\section{Introduction}
\label{sec:intro}

This paper forecasts thirty-minute-ahead realized variance for the
S\&P~500 complex on a near-twenty-four-hour panel running from January
1998 to April 2024 --- over a quarter million scored half-hour bars,
most of them outside regular trading hours. The question is not whether
volatility is predictable; on this panel the target's own lag ladder
carries the overwhelming share of the predictable variation; the lag
ladder is the benchmark throughout. The questions are what, if anything,
forty-one exogenous channels --- higher-order moments, liquidity,
cross-sectional aggregates, sentiment, implied volatility, and options
order flow --- add to that backbone; which estimator class can extract
it; and what structure the extractable nonlinearity has. Every answer is
measured out of sample, bar by bar, under one common evaluation
protocol, and every claim is paired against the same per-bar
least-squares benchmark (QLIKE \unifAzeroQLIKE{}).

Four findings organize the paper. The first is about the data itself.
On a mixed-frequency panel, the representation of missing and stale data
is a larger determinant of forecast accuracy than any modeling decision
we report: at any given bar, over a third of the exogenous design is
forward-filled rather than freshly printed, and a construction that
cannot distinguish the two produces single-window failures orders of
magnitude larger than every model-class effect in the paper. The
preparation pipeline --- availability indicators computed before the
fill, masked scaling statistics, neutral substitution --- is therefore
presented in full, with the incident record attached.

The second finding is about evaluation. Models are fit on a
variance-stabilized scale and scored on the raw variance scale, and the
back-transform between them is not neutral: squaring a calibrated
forecast is biased low by exactly the model's own residual variance, so
the correction is part of each arm's forecast. We replace the standard
static scalar correction with a calibrated, strictly causal retransformation protocol --- a
Mincer--Zarnowitz mean map re-estimated per window and an adaptive
second moment seeded by the previous window and updated bar by bar ---
which beats the retrospective \citet{Duan1983} convention it replaces on
every arm measured, and beats it causally. A mean-reverting
GARCH$(1,1)$ refinement \citep{Bollerslev1986,EngleMezrich1996} of
the same recursion --- the covariance-stationary case, not the
integrated EWMA of \citet{EngleBollerslev1986} --- improves further
still and is reported as an available, not adopted, upgrade.

The third finding is about the exogenous information and the estimator
jointly --- they cannot be separated. The exogenous channels carry
real, statistically significant information, but none of it is
recoverable by an unshrunk fit at this width: selection and unpenalized
joint fits fail, and only shrinkage recovers the signal. The head-to-head that
follows locates the structure precisely: ridge beats lasso on every
design, the lasso's zeroes fall where zeroing is free, and no small set
of columns or principal directions concentrates the signal. The
coefficient structure is \emph{dense but weak}, the regime of
\citet{GiannoneLenzaPrimiceri2021} and \citet{ShenXiu2025} --- many nonzero
coefficients, none large --- and the estimator built for it is ridge
with a block-diagonal Tikhonov penalty: the
strong low-dimensional backbone barely penalized, the wide exogenous
block shrunk two orders of magnitude harder, with fixed penalties and
no selection step anywhere.

The fourth finding is about nonlinearity. Gradient-boosted trees beat
the best linear arm on all nine feature sets, but by a margin an order
of magnitude smaller than the information effect, and decompositions of
the tree fits attribute the bulk of that margin to additive-plus-pairwise
structure --- products of columns the panel already contains, gated on
the session clock.
% PRODUCT BLOCK PARKED 2026-08-18: the product-block/three-block conclusion of the fourth finding is held out of the rendered paper; the paragraph now closes on the two-block ridge.
\iffalse
Writing those products down as one hundred frozen
columns and shrinking them as a third block converts the tree's
specification into the next block of the same Tikhonov penalty:
one solver, three fixed penalties, refit exactly at
every bar, scoring \unifBlkThreeUserQLIKE{} against the benchmark's
\unifAzeroQLIKE{} (paired increment \unifIncrBlkThreeUserDQ{} at DM
\unifIncrBlkThreeUserDM{} for the product block alone).
\fi
That structure is a specification statement, not yet an estimator: the
paper's linear ladder ends at the two-block ridge --- one solver, two
fixed penalties, refit exactly at every bar --- which scores
\unifBlkTwoUserQLIKE{} at DM \unifBlkTwoUserDM{} against the
benchmark's \unifAzeroQLIKE{}. The paper closes by reading the
forecast against the same-day index-option market: the option-implied
variance is informationally redundant given the model, and a
delta-hedged near-the-money book on the model's signal clears quoted
spreads (Section~\ref{sec:spxw_econ}).

A methodological theme runs through all four: discipline beats tuning.
Every hyperparameter in the campaign is either a fixed constant or
selected causally on a short forward validation tail, and the campaign
measures what that selection buys --- repeatedly, less than it costs.
Tuned elastic nets lose to tuned ridge not because the data prefer dense
coefficients but because penalty mistakes are cheap for ridge and
destructive for selectors; a fixed forgetting half-life helps and a
tuned one eliminates the gain; fixed block penalties outperform their
causally tuned counterparts. The final model's constants are justified
individually by measurement and never optimized on evaluation loss.

\emph{Roadmap.} Section~\ref{sec:related} places the paper in the
literature. Section~\ref{sec:data_prep} defines the panel, the
preparation pipeline, and the exogenous buckets.
Section~\ref{sec:methods} fixes the evaluation framework --- the
benchmark, the testing machinery, and the back-transformation protocol ---
and then the methods: the estimators and
the block constructions. Section~\ref{sec:results} reports the
results in the order summarized above.
Section~\ref{sec:conclusion} concludes. The appendix carries the full
arm-level master table and the protocol notes.
